\chapter{Manual de software}

Este capítulo recoge los extractos de código más relevantes y las
instrucciones de uso del framework. El código fuente completo está disponible
en los repositorios del proyecto: el backend \emph{uniback} y el frontend
\emph{gui-components} (librería \texttt{ngt-gui}).

\section{Despliegue}
El sistema completo se levanta con un único comando gracias a la
configuración \texttt{docker compose} descrita en el capítulo de Desarrollo
(Algoritmos~\ref{lst:dockercompose} y~\ref{lst:dockerfile}):

\begin{lstlisting}[style=docker, caption={Levantar el entorno completo}, label={lst:deploy}]
docker compose up --build
# API:        http://localhost:8000
# PostgreSQL: localhost:5433  (db ub_db, user postgres)
\end{lstlisting}

\section{Uso del backend como librería}
Un proyecto consumidor inicializa el framework con un diccionario de
configuración y añade sus propios modelos sobre la base ORM devuelta, sin
modificar el núcleo.

\begin{lstlisting}[language=Python, caption={Inicialización y modelo propio del consumidor}, label={lst:init}]
from uniback import initialize

config = {
  "database": {"url": "postgresql://postgres:pg_passwd@ub_pg:5432/ub_db"},
  "api": {"title": "Mi API", "version": "1.0.0"},
}
orm_base, session_factory, app = initialize(config)

from sqlalchemy.orm import Mapped, mapped_column
from sqlalchemy import String

class MiEntidad(orm_base):
    __tablename__ = "mi_entidad"
    id: Mapped[int] = mapped_column(primary_key=True)
    name: Mapped[str] = mapped_column(String(100))
\end{lstlisting}

\section{Exposición de una entidad y consumo del contrato}
Tras exponer la entidad con una fábrica CRUD, queda disponible toda la
convención de endpoints (listar, obtener, crear, actualizar, borrar y
\texttt{schema.json}). El Algoritmo~\ref{lst:curl} muestra el ciclo de prueba
rápido con \gls{curl} usando el \emph{login} local de desarrollo.

\begin{lstlisting}[style=docker, caption={Prueba rápida del contrato con curl}, label={lst:curl}]
# Login local de desarrollo (solo dev)
curl -c cookies.txt -X PUT "http://localhost:8000/api/authn?user=admin"
# Esquema enriquecido de la entidad
curl -b cookies.txt http://localhost:8000/roles/schema.json
# Listado (envoltorio: content / count / issues)
curl -b cookies.txt http://localhost:8000/roles/
# Bundle versionado de todos los esquemas
curl -b cookies.txt -i http://localhost:8000/api/sys/schemas
\end{lstlisting}

\section{Generación del formulario en el frontend}
En la aplicación Angular, registrar el módulo de tipos Formly y declarar el
componente dinámico es suficiente para obtener un formulario completo a partir
del \emph{path} de la entidad.

\begin{lstlisting}[language=Java, caption={Formulario dinámico en una vista Angular}, label={lst:front}]
// app.module: cargar FormlyComponentsModule (tipos ng-zorro + custom)
@Component({
  template: `
    <ngt-dynamic-form entityPath="/roles"
                      [(model)]="model" mode="create"
                      (submitted)="save($event)">
    </ngt-dynamic-form>`
})
export class RoleCreatePage {
  model: Record<string, unknown> = {};
  save(value: unknown) { /* EntityClient.create() ya invocado */ }
}
\end{lstlisting}

\section{Cómo escribir un plugin}
La forma recomendada de añadir un dominio nuevo no es tocar el núcleo, sino
escribir un \gls{plugin} (sección~\ref{sec:plugins}). Un plugin es un paquete
Python dentro de \texttt{src/plugins/}; el \texttt{PluginManager} lo descubre
solo al arrancar buscando \texttt{<paquete>/plugin.py} y, dentro, una subclase
de \texttt{UnibackPlugin}. No hay que registrarlo en ningún sitio. La estructura
mínima es:

\begin{lstlisting}[style=docker, caption={Estructura mínima de un plugin}, label={lst:plugin-tree}]
 src/plugins/plants/
 |-- __init__.py   # docstring del paquete
 |-- models.py     # modelo SQLAlchemy del dominio
 `-- plugin.py     # la clase UnibackPlugin con los hooks
\end{lstlisting}

El modelo hereda de \texttt{FunctionalObject} para reaprovechar
\texttt{id}/\texttt{uuid}/\texttt{name}, propiedad automática, marcas de tiempo,
borrado lógico y polimorfismo. El tipo de cada columna SQLAlchemy determina el
\emph{widget} del formulario (\texttt{Date} $\rightarrow$ \emph{datepicker},
\texttt{Enum} $\rightarrow$ \emph{select}, \texttt{Float} $\rightarrow$
numérico\dots), y \texttt{info=\{"ui": \dots\}} permite forzar un \emph{override}
manual.

\begin{lstlisting}[language=Python, caption={Modelo de dominio de un plugin (\texttt{plants/models.py})}, label={lst:plugin-model}]
class Plant(FunctionalObject):
    __tablename__ = f"{get_table_prefix()}plants_plants"
    __mapper_args__ = {"polymorphic_identity": data_object_type_id["plant"]}
    id: Mapped[int] = mapped_column(
        BigInteger,
        ForeignKey(f"{get_table_prefix()}functional_objects.id"),
        primary_key=True,
    )
    species: Mapped[str | None]   = mapped_column(String(200))
    latitude: Mapped[float | None]  = mapped_column(Float)   # -> input numerico
    longitude: Mapped[float | None] = mapped_column(Float)
    planted_on: Mapped[date | None] = mapped_column(Date)    # -> datepicker
    health: Mapped[PlantHealth | None] = mapped_column(      # Enum -> select
        SAEnum(PlantHealth, name="plant_health"))
    notes: Mapped[str | None] = mapped_column(
        String(2000),
        info={"ui": {"ui-widget": "textarea"}},              # override manual
    )
\end{lstlisting}

La clase del plugin engancha el modelo y su CRUD a los \emph{hooks}
(Algoritmo~\ref{lst:plants-plugin}). Con eso, y sin escribir nada más, la
entidad aparece en \texttt{/api/sys/entities}, se incluye en el \emph{bundle}
versionado, publica su \texttt{schema.json} y queda lista para que
\texttt{<ngt-dynamic-form>} genere su formulario en el frontend. Lo único que se
escribe a mano en la interfaz es lo propio de la app (en el ejemplo de las
plantas, el mapa). El recorrido completo está documentado en
\texttt{docs/TUTORIAL\_APP\_PLANTAS.md}.

\section{Estructura de los repositorios}
El backend sigue la estructura \texttt{src/uniback/\{persistence, services,
api, authorization, config\}} con pruebas en \texttt{tests/} y documentación
del contrato en \texttt{docs/} (\texttt{CONTRACT.md},
\texttt{DYNAMIC\_FORMS.md}, \texttt{SCHEMA\_GENERATION.md}). El frontend sigue
\texttt{projects/ngt-gui/src/lib/\{core, services, components, models\}} con la
API pública re-exportada en \texttt{public-api.ts}. Esta documentación interna
es la referencia de uso para los desarrolladores que adopten el framework.
